% Generated by scripts/generate_assets.py; do not edit.
\begin{table}[!htbp]
\centering
\normalsize
\caption{Sensibilidades condicionais com a ponte horária recalibrada}\label{tab:sensitivity}
\setlength{\tabcolsep}{4pt}
\begin{tabular}{llrrrr}
\toprule
Parâmetro & Valores avaliados & Bil. 40h & Bil. 36h & Acima 40h & Acima 36h \\
\midrule
$E_Q$ & \shortstack[l]{0,4; 0,6\\0,8; 1,0} & 1,15 a 2,40 & 7,42 a 8,85 & 1,11 a 2,40 & 6,08 a 7,42 \\
$h^\star$ & \shortstack[l]{38,0; 39,0\\40,0; 41,0} & 1,50 a 1,69 & 6,55 a 11,60 & 1,48 a 1,53 & 6,25 a 6,52 \\
$\eta_I$ & 0,33; 0,40; 0,50 & 1,57 a 1,57 & 8,38 a 8,38 & 1,53 a 1,53 & 6,52 a 6,52 \\
$\sigma$ & \shortstack[l]{0,600; 1,000\\1,326; 1,500; 2,000} & 1,39 a 1,72 & 7,50 a 9,12 & 1,36 a 1,68 & 5,82 a 7,11 \\
\bottomrule
\end{tabular}
\par\smallskip
\begin{minipage}{0.98\linewidth}\normalsize Fonte: 64 simulações com PNAD 2024T4 e base formal limitada a 44h. As quatro últimas colunas mostram o mínimo e o máximo de $A_{\rm req}$ (\%) entre os valores listados; os pontos individuais estão em \texttt{generated/sensitivity\_empirical.csv}. Cada dimensão varia separadamente, incluindo pontos interiores; as demais hipóteses permanecem na base. A cada ponto, $\omega$ e as cunhas são recalibrados ao mesmo prêmio horário e à informalidade inicial. $E_Q=1$ implica ausência de fadiga. Trata-se de uma grade de sensibilidade, sem interpretação de conjunto identificado.\end{minipage}
\end{table}
